%%% FUTURE WORK — broad brainstorm, not precise/implementation-level.
%
% Relationship to the Limitations section: Limitations (methodology chapter + the precise,
% implementation-tied Future Work items in
% documentation/methodology_critical_review_for_future_work_2026-08-15.md) covers concrete,
% actionable fixes to THIS project's actual design choices -- things like weighting RQ3's per-plot
% target by survey-history length, or a proper per-target XGBoost search. Those items answer "what
% would we do differently within the same overall approach." This document is different in kind:
% it asks whether a fundamentally different approach -- different data modality, different physics,
% different model class -- would have answered the same three research questions better. Limitations
% comes first in the dissertation and feeds into this section, but this section is broader and more
% speculative by design; not everything here needs to be feasible within a masters-dissertation
% timeline, that's expected of a brainstorm.
%%%

\section{Future work: broader methodological directions}
\label{sec:futurework-broad}

The methodology in this dissertation works entirely on tabular, plot-level summary statistics
derived from LiDAR (95th-percentile height, canopy cover, etc.), matched to point-extracted
environmental covariates. This was a deliberate, defensible scope choice for a masters
dissertation -- tabular data is fast to iterate on, every model family used here has mature,
well-understood tooling, and the resulting attribution methods (coefficients, SHAP) are directly
interpretable. But it is worth being honest that this scope choice also sets a ceiling on what the
project can find, in ways that are more fundamental than any single hyperparameter or split
decision. The directions below describe what sits above that ceiling.

\subsection{From plot-level tabular summaries to individual-tree measurement}
\label{sec:futurework-imagery}

Every field in this dataset is a plot-level summary, not an individual-tree measurement
(\S\ref{sec:data-strengths-limitations}) -- the data cannot resolve tree-to-tree competition,
individual mortality, or within-plot structural variation, and it is also true that any measurement
or scanning error, being effectively summarised over an entire $20\text{m}\times20\text{m}$ plot,
cannot be separated from genuine biological signal at that resolution. This is not just a data
limitation to note in passing -- it is arguably the main reason a deep neural network in this
project (the DNN/CR-PINN family) never gets to do what neural networks are actually good at.
Both are fed a short vector of hand-engineered scalar summaries per plot-survey row; there is no
raw spatial structure in the input for a network to learn representations from, so the DNN's
capacity advantage over Elastic Net or XGBoost has little to draw on.

A more ambitious future direction would work from the LiDAR point cloud or canopy height model
(CHM) raster directly, rather than from percentile summaries computed over it. Individual tree
detection/delineation (ITD) methods -- watershed segmentation on the CHM, or deep-learning point
cloud segmentation architectures increasingly used in forestry remote sensing -- can recover
individual-tree height, crown extent, and local density directly from the same LiDAR acquisitions
already available for Aberfoyle. This would move the whole project's unit of analysis from a
$20\text{m}\times20\text{m}$ plot to an individual tree, which is both a more precise measurement
(no within-plot averaging over trees of different sizes/ages/competitive positions) and a more
interesting one -- it would let a future project actually ask whether an individual tree's local
neighbourhood (crowding, gaps, edge proximity) explains its growth trajectory, a question the
current plot-averaged design cannot pose at all, only gesture at via crude proxies like
\texttt{CanopyCover}.

\subsection{Visual encoders and a reconstructed forest digital twin}
\label{sec:futurework-digital-twin}

A further step beyond individual-tree detection: use a vision encoder (a CNN or vision-transformer
backbone, pretrained on canopy imagery or trained from scratch on the CHM/point-cloud data
available here) to learn a general-purpose representation of canopy structure, rather than
hand-designing the ITD pipeline's segmentation rules. Interpreting and reconstructing every
individual crown across a compartment this way effectively builds a digital twin of the forest's
current structure -- a persistent, queryable 3D/2.5D representation rather than a table of
plot-level scalars. This is a substantially larger undertaking than anything in the current
dissertation (a different data modality, a different modelling paradigm, and a much larger
engineering surface -- point-cloud/imagery pipelines, not tabular ML), and is offered here as a
longer-horizon direction rather than an immediate next step. Its main payoff would be enabling
individual-tree-level growth and competition modelling (\S\ref{sec:futurework-imagery}) at scale
across the whole study area, rather than one hand-built pipeline run once.

\subsection{Physics-based wind modelling}
\label{sec:futurework-wind-cfd}

Worth being precise about what the current wind/terrain covariates actually are, since it is easy
to assume they are already physics-based simulations when they are not. TOPEX
(\S\ref{sec:data-environmental}) is a static, purely geometric calculation from horizon angles
around each plot -- it quantifies shelter/exposure from terrain shape alone, with no fluid dynamics
involved. The Global Wind Atlas is closer to physics-based (a macro-scale wind-flow model, built
originally for wind-energy resource assessment), but it is still a coarse, $250\text{m}$-resolution
long-term climatological average -- it does not resolve local turbulence, canopy-edge acceleration,
or gap dynamics at the stand scale. Neither source is a computational fluid dynamics (CFD)
simulation of how wind actually moves through and around this specific forest's terrain and canopy
structure.

A genuine future direction is to bring in either a full local CFD simulation (e.g.\ an
OpenFOAM-style canopy-flow model resolving wind speed and turbulence intensity at the stand or
even individual-tree scale) or a mechanistic forestry wind-risk model purpose-built for this
species and region -- ForestGALES, the UK Forestry Commission's operational decision-support tool
for wind risk, which estimates critical wind speeds for stem breakage and overturning from tree
mechanics and site conditions rather than from a geometric or macro-scale proxy \ct{Gardiner,
Peltola, Kellomaki (2000), "Comparison of two models for predicting the critical wind speeds
required to damage coniferous trees," Ecological Modelling -- this is the mechanistic model
underlying GALES/ForestGALES and I'm reasonably confident it's correct, but the specific
ForestGALES \emph{tool/software} citation (commonly a Gardiner et al./Forest Research technical
publication) is a separate reference I'm not confident enough on the exact year/authors to state --
confirm both before citing}. Either would give
RQ2/RQ3's attribution models a mechanistically grounded wind covariate to test against, instead of
TOPEX's static geometry or GWA's coarse climatology -- and since this dissertation's own data
already comes from Forest Research, a ForestGALES-based covariate in particular would be a natural,
well-precedented extension rather than an exotic one.

\subsection{Foundation-model geospatial embeddings (AlphaEarth)}
\label{sec:futurework-alphaearth}

Newer satellite-derived geospatial embedding models (e.g.\ Google DeepMind's AlphaEarth
Foundations \ct{Google DeepMind, "AlphaEarth Foundations," 2025 -- I'm not confident enough on the
exact author list or publication venue (likely an arXiv preprint/technical report rather than a
peer-reviewed journal, given how recent it is) to state either -- confirm before citing}) compress
multi-modal Earth observation data into a compact per-location embedding
vector, intended for downstream similarity and classification tasks without hand-selecting
individual variables the way \S\ref{sec:feature-final}'s Set1--5 protocol does.

\textbf{Why not used here (justification, not an oversight):} the strongest reason is a genuine
methodological conflict with this dissertation's aims, not just timing. RQ2/RQ3 are explicitly
attribution tasks -- they ask \emph{which} environmental variable explains a departure from
expected growth, and in \emph{which direction}, using coefficients and SHAP values that decompose
a prediction into named, physically meaningful drivers. An embedding vector is, by design, opaque:
its dimensions do not correspond to named physical quantities, so it cannot support that kind of
attribution at all -- feeding it into RQ2b/RQ3 would trade away the entire interpretability
apparatus those research questions depend on. A secondary, less certain reason is availability:
these embedding products are recent, and it is not confirmed here whether a usable release existed
early enough in this project's timeline to have been a realistic option.

\textbf{Where it could help in future work:} not as an attribution-model input, but for tasks where
raw predictive power or similarity matters more than interpretability. Two candidates: (1) a
plot/area similarity search -- using embedding distance to find structurally comparable stands
outside Aberfoyle, which would directly address the generalisability limitation stated in
\S\ref{sec:data-strengths-limitations} (this is a single-forest proof of concept) by giving a
principled way to identify where else the findings might transfer; (2) as an additional predictive
feature for RQ1-style top-height prediction specifically, where the DNN/PINN family already
tolerates lower interpretability in exchange for accuracy, unlike RQ2b/RQ3's attribution models.

\subsection{Statistical attribution methods vs.\ neural networks -- revisiting this project's own evidence}
\label{sec:futurework-stats-vs-nn}

Worth stating precisely, because the intuition turns out to already be supported by this project's
own results, not just general ML folklore. Every attribution model actually used for RQ2b and RQ3
is a statistical or tree-based method with a mature, standardised interpretability toolkit: Elastic
Net (linear coefficients), XGBoost (gain importance, SHAP), the two-stage mixed-effects model
(fixed-effect coefficients), and GNNWR, which despite being a small neural network, is explicitly
constrained to be interpretable -- it multiplies a frozen, pre-fitted global OLS coefficient by a
learned local weight, rather than learning an opaque end-to-end mapping. The one genuinely
unconstrained deep architecture in the project, the DNN/CR-PINN family, is used only for RQ1's raw
predictive task, never for attribution -- and RQ2a's own finding (the residual-shrinkage pattern
motivating PINN's physics guidance) turned out to hold just as well for XGBoost, with the actual
case for continuing with PINN resting on reproducibility and its structural
$y_{\max}$/$k$ parametrisation, not on it explaining anything more clearly than a tree-based model
does.

So the more precise version of the claim is: it is not that statistical methods are
\emph{inherently} better than neural networks at attribution in general (there is active research
on interpretable/explainable deep learning -- attention-based attribution, integrated gradients,
concept-based explanations); it is that at this project's scale (thousands of rows, roughly 30
features) and for this project's specific goal (named, defensible, direction-and-magnitude
attribution rather than raw predictive accuracy), classical statistical and tree-based methods are
better matched, and their attribution tooling is more mature and easier to defend rigorously than
the equivalent tooling for deep networks. That is a scale-and-purpose-specific claim, evidenced by
this project's own design choices and RQ2a's finding, not a general law about neural networks --
worth stating it that way rather than overclaiming.

\subsection{A better method for RQ3, independent of AI/ML: hierarchical, jointly-fitted models}
\label{sec:futurework-hierarchical}

RQ3's central difficulty is not model choice, it is that each plot contributes only a handful of
repeated observations to its own asymptote estimate -- a small-$N$-per-unit problem, not a
wrong-model-family problem. The textbook answer to that specific problem needs no AI/ML at all: a
\textbf{hierarchical (multilevel) model with partial pooling}, where a plot's asymptote estimate is
shrunk toward its compartment or yield-class expectation in proportion to how little data that
plot actually has, rather than every plot's point estimate being trusted equally regardless of
whether it came from two surveys or six. Concretely, this looks like a Bayesian hierarchical
growth-curve model (e.g.\ fitted with Stan/\texttt{brms}), with plot-level random effects on
$y_{\max}$ nested inside compartment, and environmental covariates entering as fixed or
random-slope predictors, fit \emph{jointly} in one stage rather than the current design's two
disconnected stages (Stage 1: closed-form per-plot asymptote fit, its own uncertainty then
discarded; Stage 2: attribution model treating that point estimate as measured without error,
\S\ref{sec:rq3-target}).

This is not a new idea introduced here -- it is, precisely, the "full non-linear mixed-effects
model" that RQ2b's own methodology (\S\ref{sec:rq2b}) already names as the more statistically
correct alternative to its two-stage approximation, declined there for complexity and time. The
honest framing for Future Work is therefore not "a better method exists that wasn't considered," it
is "the better method was identified and consciously deferred, for a stated and defensible reason,
and revisiting it is the single most textbook-correct upgrade available to RQ3 specifically."

\subsection{Causal design: mediators, confounders and colliders, not just correlation screening}
\label{sec:futurework-causal-design}

A worked example makes the gap concrete. Consider elevation: it plausibly affects growth deviation
through several different downstream pathways at once -- greater rainfall exposure directly, but
also, combined with slope and aspect, greater surface run-off, which in turn implies shallower
effective soil depth and moisture retention, which in turn weakens root anchorage, which in turn
increases vulnerability to windthrow -- and that whole chain is sharpened further if a drought or
storm happens to fall inside a given plot's specific survey interval. None of the intermediate
steps in that chain (soil moisture, root anchorage) are directly measured in this dataset. When an
unmeasured \emph{mediator} sits between a measured feature and the outcome, the coefficient
actually estimated on that feature (elevation, say) is a mixture of whatever direct effect exists
plus everything routed through the missing mediator -- unproblematic for pure prediction, actively
misleading if read as "elevation is the driver" rather than "elevation is a proxy for a bundle of
largely unmeasured mechanisms."

This is a different problem from what VIF screening (\S\ref{sec:feature-final}) solves. VIF only
asks whether a variable is numerically redundant with another -- a multicollinearity/stability
concern. It does not ask which of three structurally different roles a correlated variable plays:
a \textbf{confounder} (a common cause of both a variable of interest and the outcome, which must be
adjusted for to avoid bias), a \textbf{mediator} (an intermediate step on the causal path itself,
which must \emph{not} be adjusted for if the goal is the total effect), or a \textbf{collider} (a
variable caused by two other things, where conditioning on it manufactures a spurious association
between them that did not exist before). The current screening pipeline treats every correlated
variable the same way; a causally-informed design would not. \texttt{CanopyCover}'s
reverse-causation concern (\S\ref{sec:rq2b}), already investigated empirically, is arguably better
described in exactly these terms -- closer to a mediator/collider role than a simple confounder --
and revisiting it with this vocabulary would sharpen, not overturn, the existing conclusion.

Further plausible causal chains worth naming, as illustrations of the same failure mode recurring
under different variables:
\begin{itemize}
\item \textbf{Thinning has two opposing pathways}: more growing space (positive, via reduced
  competition) and more wind exposure (negative, via reduced canopy shelter). A small or null
  linear coefficient on thinning could represent two real, cancelling effects rather than "thinning
  doesn't matter."
\item \textbf{Yield class may itself be causally downstream of historical site assessment}: if
  denser planting or higher-yield-class assignment historically tracked a forester's own judgement
  of site quality, RQ3's baseline (the yield-class-implied asymptote,
  Equation~\ref{eq:local-ymax-diff}) already partly encodes the same environmental information the
  research question attributes deviations \emph{from} -- a selection/allocation confound, distinct
  from but structurally similar to the \texttt{CanopyCover} case.
\item \textbf{Edge proximity is sign-ambiguous by mechanism}: more light (positive) versus more
  wind exposure (negative) act in opposite directions. A coefficient's sign alone, without a stated
  mechanism, is correlation with a label attached, not attribution.
\item \textbf{Timing, not just exposure score, is the real confound} in the drought/storm example
  above: two plots with identical static "exposure" could show opposite deviation depending on
  whether an extreme event happened to fall inside their specific survey window -- a time-varying
  omitted variable that a \emph{persistent} target (\S\ref{sec:rq2b}, \S\ref{sec:rq3-target}) is
  structurally unable to represent, connecting directly to
  \S\ref{sec:futurework-temporal-spatial} below.
\item \textbf{Species-scope exclusion may discard real cross-species signal}: pests/disease known
  not to affect Sitka spruce directly are excluded by the single-species scope
  (\S\ref{sec:data-source}), but if neighbouring or historically mixed stands carried
  disease/pest pressure that shaped current stand composition or management response near a plot,
  removing that information entirely is not obviously costless -- worth flagging as an open
  question rather than an assumed-safe simplification.
\end{itemize}

\subsection{Toward a causal-linked attribution model}
\label{sec:futurework-causal-model}

Two tiers, from achievable now to requiring data this project does not have:

\textbf{Achievable without new data collection.} Formally draw the assumed causal graph (a DAG:
elevation $\to$ \{rainfall exposure, slope-mediated run-off\} $\to$ soil depth/moisture $\to$ root
stability $\to$ deviation, plus a separate management/selection-confound branch through yield
class) and use it to determine the correct \emph{adjustment set} for a given causal question --
which variables to condition on to isolate a specific effect, explicitly excluding known mediators
and colliders from that set, rather than admitting every VIF-passing variable into one regression
indiscriminately. This changes which variables enter the model and why, not just how they are
screened for redundancy, and could be described and partially applied to the existing dataset
without collecting anything new.

\textbf{A middle tier that still uses ML, but targets a causal rather than associational estimate.}
Double machine learning / causal forests \citep{chernozhukov2018double} use flexible ML models
(potentially including the tree ensembles already used throughout this project) to model nuisance
relationships (treatment assignment, outcome surface) while still targeting a causally-interpretable
effect estimate with valid uncertainty -- a genuinely different exercise from reading a raw SHAP
value as if it were a causal effect, and one that keeps this project's existing tooling relevant
rather than discarding it. \ct{Chernozhukov, Chetverikov, Demirer, Duflo, Hansen, Newey, Robins
(2018), "Double/debiased machine learning for treatment and structural parameters" -- full
reference (journal/volume/pages, BibTeX key) not yet confirmed, find and add}

\textbf{The strong version (instrumental variables, natural experiments), needing data not
currently available.} Requires a source of variation in an environmental or management driver that
is exogenous -- uncorrelated with unmeasured confounders -- such as a historical policy change or a
sharp, arbitrary management threshold. No obvious candidate is identified in the current dataset;
worth naming as the gold standard while being explicit that pursuing it would likely require new
historical records (e.g.\ archived planting/management policy decisions) rather than a change of
model alone.

\subsection{Spatial and temporal: separate models, with one shared prerequisite}
\label{sec:futurework-temporal-spatial}

The principled answer -- a single joint spatiotemporal model, with spatial random effects,
temporal random effects, and a spatial$\times$temporal interaction term, capable of representing
exactly the "drought hits exposed slopes harder" pattern described above -- is not recommended for
this project as it currently stands. That kind of model needs temporal replication this dataset
does not have: a handful of irregular, multi-year survey gaps per plot is far too little to fit a
spatiotemporal covariance structure without the temporal component being dominated by noise. Two
separate, simpler models -- one for "does location matter" (\S\ref{sec:rq2b}, \S\ref{sec:rq3-target}),
one for "does time matter" -- remain the more defensible choice given the data actually available.

That does not mean the spatial$\times$temporal interaction idea is a dead end, only that it should
be tested cheaply rather than via a full joint model. It already was tested, in exactly this form:
\texttt{temporal\_wind\_extension\_check.py} interacted a cohort-level storm profile against
plot-level static exposure (\S\ref{sec:futurework-broad}'s discussion in the companion review
document, \texttt{documentation/methodology\_critical\_review\_for\_future\_work\_2026-08-15.md})
and found no gain. As established there, that result is best read as a target-shape problem, not a
refutation of the interaction idea itself: a persistent per-plot target structurally cannot reward
interval-specific signal, interaction terms or not. The conclusion from that check and the
conclusion here are the same fix, reached from two different directions -- an interval-level target
(survey-to-survey deviation \emph{change}, naturally extending RQ2a's row-level shape rather than
RQ2b/RQ3's persistent one) is the shared prerequisite for testing \emph{both} the temporal question
on its own \emph{and} any spatial$\times$temporal interaction meaningfully. Build that target once,
and both open questions become testable from the same design, rather than needing two separate
future extensions.


--- # Critical review of the implemented methodology (either gores into limitations for methodlogy or if directly connected to results then their, else future work) — prep for Future Work

**Purpose**: not dissertation prose. A working audit of every real decision made in the actual
pipeline (data cleaning through each RQ's experimental setup), each with its justification, its
honest limitation, and why it was kept anyway — then a consolidated list of what would be *more
obviously done differently* if starting over, for the Future Work section specifically.

**Sourcing**: data-cleaning facts are drawn from `documentation/august_draft/3_Chapter_data/
data_chapter_2026-08-13_1222.tex` (authoritative for what was actually applied). A few
exploratory data-quality observations from earlier in the project (LAI 2021 anomaly, `plyr=0`
artefact) are noted separately as *observed, not necessarily gated by a formal filter* — flagged
where that distinction matters, since I haven't re-read the cleaning notebook line-by-line in this
session. Methodology facts are drawn from the current chapter draft
(`4_Chapter_methodology/g_methodology_2026-08-15_clean_all_applied.tex`) and
`TEMP_results/README.md` / `documentation/research_questions_overview.md` for what's actually been
checked.

---

## 1. Data cleaning stage

| Decision | Justification | Limitation | Why kept anyway |
|---|---|---|---|
| Two cohorts (4survey/6survey) instead of one, split at the 2006→2008 LiDAR coverage jump | A single cohort would either discard the pre-2008 surveys' longer trajectories or discard the +37km² of post-2008 spatial coverage; splitting keeps both | 6survey is a *nested* subset of 4survey, not an independent replication — findings that hold on both aren't two separate confirmations, and 6survey's 47 compartments are too few for some attribution splits (RQ2b excludes it entirely) | Spatial coverage matters more for this dissertation's core spatial-generalisation RQs than trajectory length; 6survey still earns its keep as a temporal-split cohort and a sensitivity check |
| Restricted to Sitka spruce (SS) only | Other species (esp. broadleaves) have seasonally variable canopy cover, which would confound `CanopyCover` as a stand-structure signal; SS is also the dominant species in Aberfoyle | Every finding is species- and site-specific; nothing generalises beyond SS/Aberfoyle by design | Explicitly framed as a proof-of-concept, not a general model — a defensible scope choice, not an oversight, but worth restating plainly in Future Work rather than only in the data chapter |
| Age/height cleaning bounds (planting year post-1850 and pre-survey-year; age 0–200; height 0–60m, the local SS ecological max) | Standard biological plausibility gates | Removed ~0 rows for SS specifically — almost all attrition instead comes from requiring complete, uninterrupted LiDAR coverage | These bounds were sanity checks, not real filters for this species — fine, but means the *actual* main source of sample loss (coverage-completeness) isn't a considered methodological choice so much as a hard data constraint |
| **Age-floor filter (≥30 years by 2023) investigated but deliberately not applied** | Quantified precisely: would remove 19% of 4survey plots (meaningful) vs. only 0.9% of 6survey — a real, considered trade-off, not an oversight | Early-age rows may include less reliable LiDAR-derived heights (the CR curve's near-linear early segment, `Figure~\ref{fig:age-linearity}`, is itself partly a hedge against this), but no formal floor was ever imposed | The 19% cut on 4survey was judged too costly for the spatial-coverage goal; deferred rather than resolved — **this is one of the strongest, most concrete Future Work candidates**: revisit with an explicit, stated age-floor rule now that the RQ1 baseline/PINN comparison and the linear-regression age-range argument both implicitly lean on this same early-age reliability question |
| Target = 95th-percentile LiDAR height (`elev_percentile_95th`), used directly without correction | Matches Forest Research's own basis for yield-class/volume calculations, keeping this project's target comparable to their convention | Not independently validated against ground-truth field height for this specific dataset (accepted as the industry-standard proxy) | Standard practice, and no ground-truth field survey exists to check against anyway |
| Environmental variables missing → **plot dropped entirely**, not imputed | Avoids introducing synthetic values that could masquerade as a real environmental driver in RQ2/RQ3's attribution — imputation bias would be a serious threat to an attribution task specifically | 430 plots lost from 4survey's candidate pool, concentrated (SoilGrids: 413, CEH TWI: 39, clustered in one compartment) — this is a **non-random dropout**, so the modelled population is subtly geographically biased away from wherever that compartment's plots sit | Imputation would need its own justification and validation (a whole sub-method), which wasn't judged worth the complexity for an attribution-focused chapter where introducing any synthetic signal is riskier than losing coverage |
| Environmental variables assigned by **point extraction** at plot centroid (MAUP heuristic for clipped plots), not interpolation or zonal/area-weighted aggregation | Simpler, avoids compounding one MAUP-style choice (centroid estimation for clipped plots) with a second one (aggregation method) | Several source rasters are coarser than the 20×20m plot footprint (SoilGrids 250m, CHELSA/HadUK 1km) — many nearby plots share an identical raw value, reducing effective variance for VIF/importance ranking on those variables specifically | Most sources are already *someone else's* interpolated/modelled product upstream (SoilGrids, CHELSA) — a fancier local aggregation wouldn't recover real sub-cell variation that was never captured by the source in the first place, so the extra complexity buys little |
| `GYCspec95`/`yldc` excluded from all models | Both are calculated *from* the height target — textbook target leakage | None beyond the obvious (loses a variable that would otherwise be informative) | Non-negotiable; keeping them would invalidate every downstream result |

**Data-quality caveats observed but not necessarily formally gated** (flagged for verification, not asserted as implemented filters): an early LAI-2021 anomaly and a `plyr=0` artefact were identified during exploratory data work earlier in the project. If these were resolved by a specific filter in the final cleaning notebook, that should be stated explicitly in the data chapter rather than left implicit — worth a quick check before the dissertation is finalised, since a marker who reads the data chapter closely could reasonably ask "what happened to X anomaly you must have found."

---

## 2. Common experimental design (splits, leakage control)

| Decision | Justification | Limitation | Why kept anyway |
|---|---|---|---|
| Spatial-block (compartment-level) as the primary split | Plot-level splitting leaks location (nearest-neighbour argument, now visualised in `Figure~\ref{fig:splits-combined}`) | — | Directly validated: the plot-level integrity check found a real, model-dependent inflation gap (DNN much more than PINN) |
| 60m buffer, set from **typical plot spacing**, not a measured spatial-autocorrelation range | Simple, defensible starting heuristic | Circular to fix with a semivariogram *after* a split already exists — flagged explicitly in the chapter's own limitations | **Most obvious Future Work fix**: a two-stage procedure — fit a rough semivariogram on an initial coarse split, use its range to set the buffer, then refit the real splits — would resolve the circularity without much extra cost, and wasn't done only because of time, not principle |
| Feature-set ranking (Set1–5) done **once globally**, not per fold | Keeps every fold comparing the same feature set, which the whole model-family-comparison design depends on | Selection variance is never separated from model variance | A fold-specific ranking is the "more obvious" textbook-correct approach, explicitly traded away for comparability — worth naming directly in Future Work as a deliberate simplification, not an oversight |
| One XGBoost hyperparameter config reused across RQ1/RQ2/RQ3 | Keeps the three research questions' XGBoost models comparable to each other | Never derived from a systematic search; originated from one earlier exploratory fit and was applied unevenly (RQ1: still on library defaults in production, only sensitivity-checked; RQ2b: fixed in production; RQ3: used the shared config from the start but its own origin traces to yet another unvalidated exploratory script) | This is the one area where the honest answer is "this should have been a proper per-target hyperparameter search (e.g. Optuna, nested CV) from the outset" — not a principled trade-off, closer to technical debt that was caught late. **Say this plainly in Future Work rather than dressing it up as a considered choice** |

---

## 3. RQ1 — top-height prediction

| Decision | Justification | Limitation | Why kept anyway |
|---|---|---|---|
| Physics-loss weight fixed a priori at λ=1.0 (symmetric), not validation-selected | Avoids circular evaluation — picking λ to maximise held-out performance, then reporting that performance, would overstate the physics constraint's value | The physics-weight ablation later found λ=0 beats λ=1 on accuracy — the pre-specified default is *not* the best-performing one | This is a genuine, defensible tension worth stating directly: the project treated λ as the **object of study** (does physics guidance help, tested cleanly), not as a hyperparameter to be optimised. Most PINN papers do the latter; this one deliberately didn't, at the cost of the "guided" model's headline accuracy number |
| Shared architecture size (3×128) across all DNN/CR-PINN variants, size fixed rather than tuned per variant | Isolates the loss-formulation comparison from confounding architecture-capacity differences | — | Validated directly by the architecture-size sensitivity sweep (found robust, low sensitivity) — one of the better-supported design choices in the chapter |
| Environmental features are **static per plot** for RQ2b/RQ3's candidate pool (`plot_environmental_features.parquet`, one row per plot, zero year-varying columns) | Matches RQ2b/RQ3's own targets, which are themselves persistent per-plot summaries (`mean_cr_residual`, `local_y_max_difference`), not per-survey values | A genuinely time-varying source (MIDAS wind, interval-resolved) **was tried against RQ3's target and tested honestly, not assumed to fail** — see below — and the negative result is a real, checked finding, not an untested design gap | **Correcting my earlier claim here**: this was NOT simply left untried. `models/growth_curve_attribution/temporal_wind_extension_check.py` (run 2026-08-04, `documentation/experiment_log.md` lines 2731-2775) built real survey-interval-level MIDAS storm metrics (one row per actual survey gap, e.g. 2012$\to$2021, per station, per gust threshold, annualised) and tested them against RQ3's target under the same 5-fold spatial CV. Because that target is a single persistent number per plot, the interval metrics first had to be collapsed to one value per *cohort* (inverse-distance + interval-length weighted across stations/intervals) to attach to it at all -- and a cohort-constant number can't explain cross-*plot* variation on its own, so the check multiplied it against plot-varying exposure (`topex`, `windward_topex`, `whcl`, `gwa_wind_speed_50m`) to create 24 interaction columns. Result: essentially no gain on 4survey (0.1302$\to$0.1304 EN; XGBoost unchanged at 0.1222), a small *drop* on 6survey (0.0437$\to$0.0277 EN; XGBoost unchanged). The experiment log's own conclusion is precise: this doesn't mean storm timing is biologically irrelevant, it means a *persistent* target structurally can't reward *interval-specific* timing no matter how the join is engineered -- the fix is a differently-shaped target (e.g. survey-to-survey deviation change), not more feature engineering against the current one. **Only run for RQ3's target, not RQ2b's** -- same persistent-target shape, so the same conclusion would likely hold, but it is untested there specifically. **Scope of what this negative result actually covers**: `temporal_wind_extension_check.py` fits only Elastic Net and XGBoost (`METHODS = {"elastic_net_predicted": ..., "xgboost_predicted": ...}`) -- it never touches GNNWR, the one model in this project's toolkit built to test spatially-*varying* relationships rather than one fixed global one. So the result shows the interaction terms don't help under a globally-fixed relationship; it does not rule out the effect being real but spatially heterogeneous (e.g. terrain funnelling storm exposure differently in different parts of Aberfoyle), which is a genuinely separate, untested question, not a restatement of the one already answered. |

---

## 4. RQ2a — does environment shrink RQ1's own departure from the curve

| Decision | Justification | Limitation | Why kept anyway |
|---|---|---|---|
| No new models fitted — reuses RQ1's own predictions against the CR baseline | Cheap (seconds per combination), keeps RQ2a mechanically tied to whatever RQ1 already validated | Entirely contingent on RQ1's model roster: originally covered only DNN/PINN/PINN$_k$, and only later got an XGBoost-specific follow-up once the fairness question surfaced — not because it was designed to cover every RQ1 family from the start | The retrofit (adding XGBoost's version after the fact) worked, but it's a patch, not the original design — Future Work: build RQ2a to sweep every RQ1 model family uniformly from day one, so an upstream fairness fix doesn't require a bespoke follow-up experiment each time |

---

## 5. RQ2b — persistent departure from the shared curve

| Decision | Justification | Limitation | Why kept anyway |
|---|---|---|---|
| Age excluded from all candidate features | It's the CR baseline curve's sole predictor — including it would make the attribution target circular by construction | — | Non-negotiable, well-justified, and explicitly checked (this project's one clean example of a target-circularity audit actually catching something) |
| Two-stage mixed-effects approximation instead of a full joint non-linear mixed-effects model | Standard, citable simplification (Socha 2016); much cheaper to fit | Stage-1 uncertainty isn't propagated into Stage-2 — coefficient significance may be slightly overstated | Full joint NLME is substantially more complex to fit and wasn't judged worth it for an attribution-only (not predictive) model; a reasonable place to spend more time in Future Work if the coefficient p-values need to be defended more rigorously |
| `CanopyCover`'s reverse-causation risk checked **after** it turned out to dominate every set, not designed around from the start | The check itself (baseline-only ablation, between-compartment variance decomposition) is solid and the resulting framing (baseline control, not the headline finding) is defensible | It's a retrofit — the original design didn't anticipate needing this control | **A clear, honest Future Work item**: build the baseline-vs-environment decomposition into the attribution design from the outset for any future stand-structure variable that shares the same ALS pipeline as the target, rather than adding it reactively once a variable looks suspiciously dominant |

---

## 6. RQ3 — plot-specific growth-curve deviation

| Decision | Justification | Limitation | Why kept anyway |
|---|---|---|---|
| Asymptote-only deviation target, curve shape fixed to yield-class values, fitted by closed-form WLS | Stable with few observations per plot; directly matches what the yield-class comparison is actually about (max height potential) | Plots with different growth *rates* or inflection timing that reach the same eventual ceiling aren't treated as divergent — a real, stated scope limitation | Deliberate, well-justified scope choice, not an oversight — but worth restating in Future Work as "a companion growth-*shape* deviation metric (not just ceiling) is the natural next target" |
| Per-plot asymptote target uses whatever survey history each plot has, with **no weighting by history length** | Simpler pipeline, treats every plot symmetrically | Plots with only the cohort-minimum number of surveys get a noisier target than plots with a full history, and this noise isn't propagated or down-weighted anywhere downstream | **Probably the single most "obvious in hindsight" fix in the whole RQ3 design**: inverse-variance weighting (or even just a categorical "survey-count" control feature) by survey-history length is a standard, low-cost fix that wasn't implemented — good, concrete Future Work item |
| GNNWR kept as an attribution tool despite known 6survey fragility (sign-flips across seeds, near-zero, underpowered) | It's the only tool in the toolkit that tests spatially-*varying* coefficients at all — dropping it would remove that entire axis of RQ3 | 6survey's result there is genuinely inconclusive, not just noisy-but-directionally-right | 4survey's GNNWR result is reliable and is the headline; 6survey was always meant as a secondary check, and reporting "underpowered, no reliable direction" honestly is itself a valid finding about cohort-size requirements for spatially-varying models |

---

## Consolidated "more obvious alternative" list for Future Work

Ranked roughly by how concrete/actionable each one is, not by importance:

1. **Treat the temporal-environmental question as its own separate, purpose-built analysis with an interval-level target** (e.g. survey-to-survey deviation *change*), rather than continuing to force time-varying data into RQ2b/RQ3's persistent-target framework. This isn't a hunch — it's supported by a real, checked result: the feature-side fix was already tried (engineering cohort$\times$exposure interactions from real interval-level MIDAS data, `temporal_wind_extension_check.py`) and found not to help under two different global-relationship model families (Elastic Net, XGBoost; GNNWR was not included in that check and remains untested). The experiment log's own reasoning is the stronger argument, though: a target that collapses a plot's whole survey history into one persistent number cannot, by construction, reward interval-specific timing, regardless of how cleverly the features are engineered. That's a structural explanation, not just an empirical correlation of failure. **Design suggestion**: this future analysis would naturally extend RQ2a's shape, not RQ2b/RQ3's — RQ2a already operates at the per-survey-row level (RQ1's predictions vs. the CR baseline, matched by plot *and year*), which is inherently compatible with time-varying environmental data in a way RQ2b/RQ3's persistent per-plot summaries structurally aren't. Framing this as a proposed future RQ (rather than a buried "also try X" line) is a reasonable rhetorical choice for the Future Work section, even though structurally it's an extension of existing row-level machinery rather than a wholly new design.
2. **Inverse-variance (or survey-count) weighting for RQ3's per-plot asymptote target** — cheap, standard, directly addresses a named limitation.
3. **A real per-target XGBoost hyperparameter search** (RQ1/RQ2b/RQ3 each tuned on their own target) instead of one carried-over config applied unevenly — the most honest "this was technical debt, not a principled choice" item.
4. **Two-stage semivariogram-then-buffer procedure** to set the 60m spatial buffer empirically instead of by convention, resolving the circularity noted in the chapter's own limitations.
5. **An explicit, applied age-floor rule** for early-age LiDAR reliability — quantified during the project (19%/0.9% cut on 4survey/6survey) but never actually adopted; the linear-regression age-range argument and RQ1's baseline comparison both lean on this same unresolved question.
6. **Build RQ2a to cover every RQ1 model family from the start**, not just DNN/PINN with XGBoost bolted on later once a fairness issue surfaced.
7. **Design the baseline-vs-environment decomposition into RQ2b/RQ3 from the outset** for any stand-structure variable sharing the target's own ALS pipeline, rather than adding it reactively after `CanopyCover` turned out to dominate.
8. A companion **growth-shape** (not just asymptote) deviation target for RQ3, since the current target is explicitly scoped to ignore rate/inflection differences between plots that reach the same ceiling.

Items 1, 2, and 3 are the strongest candidates for the dissertation's actual Future Work section — each is concrete, each has a clear "why it wasn't done" (data-pipeline design, not laziness), and each maps directly onto a limitation already stated in the methodology chapter, so the Future Work section can cross-reference rather than re-derive the argument.
